\chapter{두 지도를 섞으면?}

지금 우리에게는 지도가 두 장 있어요.
허락 지도에서 놀면 엔도스랭크, 송금 지도에서 놀면 AWP예요.
그렇다면 자연스럽게 이런 생각이 들어요.
\emph{두 지도를 함께 쓰면 둘의 좋은 점을 다 가질 수 있지 않을까?}

저도 그렇게 생각했어요.
그래서 두 가지 섞는 방법을 만들어 보았어요.

\section{동전 던지기: C-PR}

교실 바닥에 두 지도를 겹쳐 놓았다고 상상해 보세요.
파란 허락 화살표와 주황 송금 화살표가 함께 그려져 있어요.

쪽지 놀이 규칙을 하나만 바꿔요.
“화살표를 따라가” 카드를 뽑으면, 동전을 한 번 더 던져요.

\begin{itemize}
\item 앞면이 나오면 \textcolor{allowcol}{\textbf{파란 허락 화살표}} 가운데 하나를 따라가요.
\item 뒷면이 나오면 \textcolor{sendcol}{\textbf{주황 송금 화살표}} 가운데 하나를 따라가요.
\end{itemize}

이 방법을 \textbf{C-PR}이라고 불러요.
C는 “짝지어 묶었다”는 뜻의 영어 낱말 coupled의 첫 글자예요.
동전이 앞면일 확률을 그리스 문자 $\lambda$(람다)로 써요.
보통 동전이면 $\lambda = 0.5$예요.
$\lambda = 1$이면 늘 허락 화살표만 따라가니 엔도스랭크와 같아지고,
$\lambda = 0$이면 늘 송금 화살표만 따라가니 AWP와 같아져요.
그러니까 C-PR은 엔도스랭크와 AWP 사이 어딘가에 있는 방법이에요.
저는 $\lambda$를 0.5로 정하고, 0.25와 0.75도 함께 시험했어요.

\begin{figure}[h]
\centering
\begin{tikzpicture}[scale=0.9, every node/.style={transform shape}]
  \node[kid] (h) at (0,0) {하준};
  \node[kid] (s) at (3.6,1.3) {서연};
  \node[robot] (r) at (3.6,-1.3) {로봇};
  \draw[allow] (h) -- node[above, sloped, font=\scriptsize, color=black] {앞면이면} (s);
  \draw[send] (h) -- node[below, sloped, font=\scriptsize, color=black] {뒷면이면} (r);
  \node[draw=sunline, fill=sunfill, circle, minimum size=9mm, font=\small] at (-1.9,0) {동전};
\end{tikzpicture}
\caption{C-PR에서는 쪽지를 넘길 때마다 동전을 던져 어느 지도의 화살표를 따라갈지 정해요.}
\end{figure}

\section{출발점 바꾸기: S-PR}

두 번째 방법은 조금 달라요.
쪽지는 송금 화살표만 따라가요.
그 대신 “선생님께” 카드가 나왔을 때 선생님이 쪽지를 아무에게나 주지 않고,
\emph{엔도스랭크 점수가 높은 친구}에게 더 자주 줘요.
믿음을 많이 받은 친구에게서 출발해서 송금 길을 걷는 셈이에요.
이 방법을 \textbf{S-PR}이라고 불러요.
S는 “씨앗을 뿌렸다”는 뜻의 seeded예요.

\section{결과를 보기 전에 정한 약속}

새 방법을 만들면 자꾸 좋은 점만 찾고 싶어져요.
그래서 저는 C-PR 점수를 하나도 계산하기 전인 2026년 9월 14일에,
C-PR이 “성공”이라고 말하려면 무엇을 넘어야 하는지 미리 적어 두었어요.
지도교수님께 보낸 계획서와 프로그램 설정 파일에 함께 적었어요.

\begin{center}
\begin{tcolorbox}[colback=white, colframe=inkblue, boxrule=0.8pt, arc=2mm, width=0.92\linewidth]
\small
\textbf{9월 14일의 약속} \quad ($\lambda = 0.5$인 C-PR에 대해)
\begin{enumerate}
\item 새 허락을 맞히는 시험에서 엔도스랭크보다 못하지 않을 것.
\item 새 송금자를 맞히는 시험에서 AWP보다 나을 것.
\end{enumerate}
두 가지를 모두 넘어야 성공이에요.
\end{tcolorbox}
\end{center}

“새 허락을 맞히는 시험”과 “새 송금자를 맞히는 시험”이 무엇인지는 13장에서 자세히 설명할게요.
지금은 이것만 기억해 주세요.
\emph{결과를 보기 전에 약속을 먼저 적었다.}
이 약속이 이 책의 끝부분에서 아주 중요해져요.

\begin{deeper}[진짜 정의]
C-PR은 점마다 두 층의 이동 확률을 섞어요. 점 $i$에서 $j$로 갈 확률은
\[
P_{ij} = \lambda\, P^{\text{허락}}_{ij} + (1-\lambda)\, P^{\text{송금}}_{ij}
\]
이고, 한 층에 화살표가 없는 점은 다른 층의 화살표를 써요.
S-PR은 송금 지도의 페이지랭크인데, 되돌림(선생님 몫)을 엔도스랭크 점수에 비례해 나누어요.
엔도스랭크 점수가 0인 점은 되돌림을 받지 않아요.
9월 14일에는 이 두 번째 규칙과 함께, 같은 기간 시험에서의 규칙(이차 규칙)도 정해 두었어요.
\end{deeper}

\begin{learned}
\begin{itemize}
\item C-PR은 쪽지를 넘길 때마다 동전을 던져 허락 지도와 송금 지도 가운데 하나를 고르는 방법이에요.
\item S-PR은 엔도스랭크가 좋아하는 친구에게서 출발해 송금 지도를 걷는 방법이에요.
\item C-PR이 성공했다고 말하는 기준은 결과를 보기 전인 9월 14일에 미리 적어 두었어요.
\end{itemize}
\end{learned}
